\documentclass{article}
\usepackage{amsmath,amssymb,amsthm}
\usepackage{enumitem}
\usepackage{hyperref}
\usepackage{geometry}
\geometry{margin=1in}
\newtheorem{theorem}{Theorem}[section]
\newtheorem{lemma}{Lemma}[section]
\newtheorem{assumption}{Assumption}[section]

\begin{document}

\section*{Algorithm Name}
\textbf{Convex Trust‑Region Method (CTR)}  

The method solves, at each iteration $k$, a quadratic model of a convex $L$‑smooth objective
\[
f:\mathbb{R}^{n}\rightarrow\mathbb{R},
\qquad 
\text{with } \|\nabla f(x)-\nabla f(y)\|\le L\|x-y\|,\;\forall x,y .
\]
A trust–region radius $\Delta_k>0$ restricts the step $p_k$:
\[
\|p_k\|\le \Delta_k .
\]

--------------------------------------------------------------------

\section*{Mathematical Formulation}
\begin{enumerate}[label=\arabic*)]
\item \textbf{Model.}  At the current iterate $x_k$ define the (exact) quadratic model
\[
m_k(p)\;:=\; f(x_k)+\nabla f(x_k)^{\top}p+\frac{L}{2}\|p\|^{2},
\qquad p\in\mathbb{R}^{n}.
\tag{1}
\]
Because $f$ is $L$‑smooth, the model satisfies the \emph{upper‑bound property}
\[
f(x_k+p)\;\le\; m_k(p), \qquad\forall p\in\mathbb{R}^{n}.
\tag{2}
\]

\item \textbf{Subproblem.}  Compute an \emph{approximate} solution of
\[
p_k \;:=\; \arg\min_{\|p\|\le\Delta_k} m_k(p).
\tag{3}
\]
Since $m_k$ is strictly convex, the exact solution is available in closed form:
\[
p_k=
\begin{cases}
-\dfrac{1}{L}\nabla f(x_k), & \displaystyle\Bigl\|\dfrac{1}{L}\nabla f(x_k)\Bigr\|\le\Delta_k,\\[2ex]
-\displaystyle\Delta_k\,
\dfrac{\nabla f(x_k)}{\|\nabla f(x_k)\|}, & \text{otherwise}.
\end{cases}
\tag{4}
\]

\item \textbf{Acceptance test.}  Let
\[
\text{pred}_k\;:=\; m_k(0)-m_k(p_k)
= -\nabla f(x_k)^{\top}p_k-\frac{L}{2}\|p_k\|^{2},
\]
\[
\text{act}_k\;:=\; f(x_k)-f(x_k+p_k).
\]
Define the ratio
\[
\rho_k\;:=\;\frac{\text{act}_k}{\text{pred}_k}.
\tag{5}
\]

\item \textbf{Update of the iterate.}
\[
x_{k+1}= 
\begin{cases}
x_k+p_k, & \rho_k\ge\eta_{1},\\
x_k,    & \rho_k<\eta_{1},
\end{cases}
\tag{6}
\]
with $0<\eta_{1}<\eta_{2}<1$ (e.g. $\eta_{1}=0.1$, $\eta_{2}=0.75$).

\item \textbf{Update of the radius.}
\[
\Delta_{k+1}= 
\begin{cases}
\min\{\gamma_{\text{inc}}\Delta_k,\;\Delta_{\max}\}, & \rho_k\ge\eta_{2},\\[1ex]
\Delta_k,                                         & \eta_{1}\le\rho_k<\eta_{2},\\[1ex]
\gamma_{\text{dec}}\Delta_k,                      & \rho_k<\eta_{1},
\end{cases}
\tag{7}
\]
where $\gamma_{\text{inc}}>1$, $0<\gamma_{\text{dec}}<1$.

\item \textbf{Termination.}  Stop when $\|\nabla f(x_k)\|\le\varepsilon$ or a maximal
iteration count $K_{\max}$ is reached.
\end{enumerate}

--------------------------------------------------------------------

\section*{Pseudocode}
\begin{verbatim}
Input:  x0, Δ0>0, Δmax>Δ0, η1, η2, γ_inc>1, γ_dec∈(0,1), ε>0, Kmax
Set k ← 0
while k < Kmax and ‖∇f(xk)‖ > ε do
    gk ← ∇f(xk)
    %--- solve the trust‑region subproblem (4) -----------------
    if ‖gk‖ ≤ L·Δk then
        pk ← -(1/L) * gk                % interior solution
    else
        pk ← -Δk * gk / ‖gk‖            % boundary Cauchy point
    end if
    %--- evaluate actual and predicted reduction ---------------
    act ← f(xk) - f(xk+pk)
    pred← - gkᵀ pk - (L/2)·‖pk‖²
    ρ ← act / pred
    %--- accept or reject the step ----------------------------
    if ρ ≥ η1 then
        xk+1 ← xk + pk
    else
        xk+1 ← xk
    end if
    %--- update the trust‑region radius ------------------------
    if ρ ≥ η2 then
        Δk+1 ← min(γ_inc·Δk , Δmax)
    elseif ρ < η1 then
        Δk+1 ← γ_dec·Δk
    else
        Δk+1 ← Δk
    end if
    k ← k + 1
end while
Output: xk
\end{verbatim}

--------------------------------------------------------------------

\section*{Dry Run Example}
Consider the one‑dimensional convex quadratic
\[
f(w)=(w-3)^{2},\qquad \nabla f(w)=2(w-3),\qquad L=2.
\]
Choose the hyper‑parameters
\[
w_{0}=0,\;\Delta_{0}=1,\;\Delta_{\max}=10,\;
\eta_{1}=0.1,\;\eta_{2}=0.75,\;
\gamma_{\text{inc}}=2,\;\gamma_{\text{dec}}=0.5 .
\]

\begin{center}
\begin{tabular}{c|c|c|c|c|c}
$k$ & $w_{k}$ & $g_{k}= \nabla f(w_{k})$ & $\|g_{k}\|$ & $p_{k}$ (by (4)) & $f(w_{k})$\\\hline
0 & $0$ & $-6$ & $6$ & $p_{0}= -\Delta_{0}\dfrac{g_{0}}{\|g_{0}\|}=+1$ & $9$\\
1 & $1$ & $-4$ & $4$ & $p_{1}= -\Delta_{1}\dfrac{g_{1}}{\|g_{1}\|}=+2$ (since $\Delta_{1}=2$) & $4$\\
2 & $3$ & $0$  & $0$ & $p_{2}=0$ & $0$\\
\end{tabular}
\end{center}

\textbf{Iteration 0.}  
$g_{0}=-6$, $\|g_{0}\|=6> L\Delta_{0}=2$, thus the boundary step $p_{0}=+1$ is taken.  
$f(w_{0})=9$, $f(w_{0}+p_{0})=4$, predicted reduction $=5$, actual reduction $=5$, $\rho_{0}=1\ge\eta_{1}$, so $w_{1}=1$ and $\Delta_{1}=2$.

\textbf{Iteration 1.}  
$g_{1}=-4$, $\|g_{1}\|=4\le L\Delta_{1}=4$, therefore the interior solution $p_{1}=-(1/L)g_{1}=+2$ is admissible.  
$f(w_{1})=4$, $f(w_{1}+p_{1})=0$, predicted reduction $=4$, actual reduction $=4$, $\rho_{1}=1$, accept step, $w_{2}=3$, increase radius $\Delta_{2}=4$.

\textbf{Iteration 2.}  
$g_{2}=0$, the subproblem yields $p_{2}=0$ and the algorithm terminates because $\|\nabla f(w_{2})\|=0\le\varepsilon$.

The method reaches the global minimiser $w^{*}=3$ in two successful steps.

--------------------------------------------------------------------

\section*{Assumptions}
\begin{assumption}[Convexity and Smoothness]\label{ass:convex}
The objective $f:\mathbb{R}^{n}\to\mathbb{R}$ is convex and continuously differentiable,
and its gradient is $L$‑Lipschitz:
\[
\|\nabla f(x)-\nabla f(y)\|\le L\|x-y\|,\qquad\forall x,y\in\mathbb{R}^{n}.
\]
\end{assumption}
\begin{assumption}[Bounded Level Set]\label{ass:bounded}
For the initial point $x_{0}$ the sublevel set
\[
\mathcal{L}:=\{x\;|\; f(x)\le f(x_{0})\}
\]
is bounded; consequently there exists a constant $D>0$ such that
$\|x-x^{*}\|\le D$ for all $x\in\mathcal{L}$, where $x^{*}$ denotes any global minimiser.
\end{assumption}
\begin{assumption}[Trust‑Region Parameters]\label{ass:params}
The algorithmic parameters satisfy
\[
0<\eta_{1}<\eta_{2}<1,\quad
\gamma_{\text{inc}}>1,\quad
0<\gamma_{\text{dec}}<1,\quad
0<\Delta_{\min}\le\Delta_{0}\le\Delta_{\max}<\infty .
\]
\end{assumption}

--------------------------------------------------------------------

\section*{Main Convergence Theorem}
\begin{theorem}[Global convergence and rate]\label{thm:conv}
Under Assumptions \ref{ass:convex}–\ref{ass:params}, the iterates $\{x_{k}\}$ generated by
CTR satisfy
\[
\min_{0\le i\le k}\|\nabla f(x_{i})\|
\;\le\;
\frac{2L\bigl(f(x_{0})-f^{*}\bigr)}{\Delta_{\min}\sqrt{k+1}},
\qquad\forall k\ge0,
\tag{8}
\]
where $f^{*}=\inf_{x}f(x)$. Consequently,
\[
\lim_{k\to\infty}\|\nabla f(x_{k})\|=0.
\]
If in addition $f$ is $\mu$‑strongly convex, then the iterates converge linearly:
\[
f(x_{k})-f^{*}\le\Bigl(1-\frac{\mu\Delta_{\min}}{L}\Bigr)^{k}
\bigl(f(x_{0})-f^{*}\bigr).
\tag{9}
\]
\end{theorem}

--------------------------------------------------------------------

\section*{Theoretical Properties}
\begin{enumerate}[label=\arabic*)]
\item \textbf{Monotonicity.}  Whenever $\rho_{k}\ge\eta_{1}$, the actual reduction
$\text{act}_{k}$ is non‑negative, hence $f(x_{k+1})\le f(x_{k})$.
\item \textbf{Stability w.r.t. $\Delta_{k}$.}  The radius never falls below $\Delta_{\min}>0$
by construction (Step (7)).  Therefore a lower bound on the predicted
reduction is always available.
\item \textbf{Hyper‑parameter sensitivity.}
Increasing $\gamma_{\text{inc}}$ or decreasing $\gamma_{\text{dec}}$ makes the
method more aggressive; the convergence guarantee (Theorem \ref{thm:conv})
remains valid for any choice respecting Assumption \ref{ass:params}.
\item \textbf{Comparison to baseline methods.}
For convex $L$‑smooth functions, standard gradient descent with step size $1/L$
achieves $f(\bar{x}_{T})-f^{*}=O(1/T)$.  The trust‑region method yields the
gradient norm bound (8), which implies the same $O(1/\sqrt{T})$ rate for
the optimality gap in the convex case, while providing a built‑in
mechanism to adapt the step length and to guarantee robustness on ill‑conditioned problems.
\end{enumerate}

--------------------------------------------------------------------

\section*{Discussion}
The CTR algorithm can be interpreted as a \emph{projected gradient step}
with an adaptive step length $\alpha_{k}=\min\{1/L,\;\Delta_{k}/\|\nabla f(x_{k})\|\}$.
The trust‑region acceptance test ensures that the model (2) is a reliable
upper bound for the true objective; consequently the iterates never increase
the objective value beyond a controllable factor.  The proof below
shows that, under the mild convexity and smoothness assumptions,
the method converges to a stationary point and enjoys a sublinear
rate that matches the classical worst‑case bound for first‑order methods.

--------------------------------------------------------------------

\section*{Preliminaries (Definitions and Lemmas)}
\begin{lemma}[Quadratic upper bound]\label{lem:upper}
If $\nabla f$ is $L$‑Lipschitz, then for any $x,p\in\mathbb{R}^{n}$
\[
f(x+p)\le f(x)+\nabla f(x)^{\top}p+\frac{L}{2}\|p\|^{2}.
\tag{10}
\]
\end{lemma}
\begin{proof}
By the fundamental theorem of calculus,
\[
f(x+p)-f(x)=\int_{0}^{1}\nabla f(x+tp)^{\top}p\,dt.
\]
Insert and add $\nabla f(x)$ inside the integral, use the Lipschitz bound,
and integrate to obtain (10).  ∎
\end{proof}

\begin{lemma}[Predicted reduction lower bound]\label{lem:pred}
Let $p_{k}$ be the solution of (3). Then
\[
\text{pred}_{k}
\;=\;-\nabla f(x_{k})^{\top}p_{k}-\frac{L}{2}\|p_{k}\|^{2}
\;\ge\;
\frac{1}{2L}\,\|\nabla f(x_{k})\|^{2}\;\min\!\Bigl\{1,\;
\frac{L^{2}\Delta_{k}^{2}}{\|\nabla f(x_{k})\|^{2}}\Bigr\}.
\tag{11}
\]
\end{lemma}
\begin{proof}
Two cases.

\textbf{Case 1.} $\| \frac{1}{L}\nabla f(x_{k})\|\le\Delta_{k}$
(i.e. interior solution).  Then $p_{k}=-(1/L)\nabla f(x_{k})$ and
\[
\text{pred}_{k}= \frac{1}{2L}\|\nabla f(x_{k})\|^{2},
\]
which coincides with the right‑hand side of (11) because the minimum term equals $1$.

\textbf{Case 2.} $\| \frac{1}{L}\nabla f(x_{k})\|>\Delta_{k}$
(boundary solution).  Then $p_{k}= -\Delta_{k}\frac{\nabla f(x_{k})}{\|\nabla f(x_{k})\|}$.
Compute
\[
\begin{aligned}
\text{pred}_{k}
&= -\nabla f(x_{k})^{\top}p_{k}-\frac{L}{2}\|p_{k}\|^{2}\\
&= \Delta_{k}\|\nabla f(x_{k})\|
    -\frac{L}{2}\Delta_{k}^{2}\\
&= \frac{L\Delta_{k}^{2}}{2}
   \Bigl(\frac{2}{L\Delta_{k}}\|\nabla f(x_{k})\|-1\Bigr) .
\end{aligned}
\]
Since $\|\nabla f(x_{k})\|>L\Delta_{k}$, the factor in parentheses is at least
$1$, thus
\[
\text{pred}_{k}\ge\frac{L\Delta_{k}^{2}}{2}
= \frac{1}{2L}\|\nabla f(x_{k})\|^{2}\,
  \frac{L^{2}\Delta_{k}^{2}}{\|\nabla f(x_{k})\|^{2}} .
\]
The right‑hand side is exactly the second argument of the $\min$ in (11).
∎
\end{proof}

--------------------------------------------------------------------

\section*{Main Proof (Step‑by‑Step Derivation)}
We prove Theorem \ref{thm:conv}.  Throughout the proof we keep the notation
\[
\text{act}_{k}=f(x_{k})-f(x_{k}+p_{k}),\qquad
\text{pred}_{k}= -\nabla f(x_{k})^{\top}p_{k}-\frac{L}{2}\|p_{k}\|^{2}.
\]

\begin{enumerate}[label=\textbf{[Step \arabic*]}]
\item \textbf{Model upper bound.}  
By Lemma \ref{lem:upper} with $x=x_{k}$ and $p=p_{k}$,
\[
f(x_{k}+p_{k})\le m_{k}(p_{k})=f(x_{k})-\text{pred}_{k}.
\tag{12}
\]
Hence
\[
\text{act}_{k}=f(x_{k})-f(x_{k}+p_{k})\ge\text{pred}_{k}.
\tag{13}
\]

\item \textbf{Relation between $\rho_{k}$ and reduction.}  
From definition (5) and inequality (13) we obtain
\[
\rho_{k}
=\frac{\text{act}_{k}}{\text{pred}_{k}}
\ge 1.
\tag{14}
\]
Thus whenever the step is accepted ($\rho_{k}\ge\eta_{1}$) we have
$\text{act}_{k}\ge\eta_{1}\,\text{pred}_{k}$.

\item \textbf{Descent inequality.}  
If $\rho_{k}\ge\eta_{1}$, then using (14) and (13):
\[
f(x_{k+1})\le f(x_{k})-\eta_{1}\,\text{pred}_{k}.
\tag{15}
\]

\item \textbf{Lower bound on predicted reduction.}  
Apply Lemma \ref{lem:pred} to obtain
\[
\text{pred}_{k}\ge
\frac{1}{2L}\,\|\nabla f(x_{k})\|^{2}\;
\min\Bigl\{1,\frac{L^{2}\Delta_{k}^{2}}{\|\nabla f(x_{k})\|^{2}}\Bigr\}.
\tag{16}
\]

\item \textbf{Two regimes.}  
Define the set
\[
\mathcal{I}:=\Bigl\{k\;\big|\; \|\nabla f(x_{k})\|\le L\Delta_{k}\Bigr\}.
\]
If $k\in\mathcal{I}$, then the minimum in (16) equals $1$ and
\[
\text{pred}_{k}\ge\frac{1}{2L}\|\nabla f(x_{k})\|^{2}.
\tag{17}
\]
If $k\notin\mathcal{I}$, the minimum equals
$\frac{L^{2}\Delta_{k}^{2}}{\|\nabla f(x_{k})\|^{2}}$ and
\[
\text{pred}_{k}\ge\frac{L\Delta_{k}^{2}}{2}.
\tag{18}
\]

\item \textbf{Bounding the number of ``large‑gradient'' iterations.}  
Whenever $k\notin\mathcal{I}$, inequality (15) together with (18) yields
\[
f(x_{k+1})\le f(x_{k})-\eta_{1}\frac{L\Delta_{k}^{2}}{2}
\le f(x_{k})-\eta_{1}\frac{L\Delta_{\min}^{2}}{2},
\tag{19}
\]
because $\Delta_{k}\ge\Delta_{\min}$ by Assumption \ref{ass:params}.
Summing (19) over all such indices gives a finite bound:
\[
\#\{k\notin\mathcal{I}\}\;\le\;
\frac{2\bigl(f(x_{0})-f^{*}\bigr)}{\eta_{1}L\Delta_{\min}^{2}}.
\tag{20}
\]

\item \textbf{Bounding the gradient norm on ``small‑gradient'' iterations.}
For $k\in\mathcal{I}$ we combine (15) and (17):
\[
f(x_{k+1})\le f(x_{k})-\eta_{1}\frac{1}{2L}\|\nabla f(x_{k})\|^{2}.
\tag{21}
\]
Rearranging,
\[
\|\nabla f(x_{k})\|^{2}
\le \frac{2L}{\eta_{1}}\bigl(f(x_{k})-f(x_{k+1})\bigr).
\tag{22}
\]

\item \textbf{Summation over all iterations.}  
Sum (22) over all $k$ that belong to $\mathcal{I}$ and use telescoping:
\[
\sum_{k\in\mathcal{I}}\|\nabla f(x_{k})\|^{2}
\le \frac{2L}{\eta_{1}}\bigl(f(x_{0})-f^{*}\bigr).
\tag{23}
\]

\item \textbf{Deriving the $1/\sqrt{K}$ bound.}  
Let $K$ be the total number of iterations.  Split the set $\{0,\dots,K\}$
into $\mathcal{I}$ and its complement.
From (23) and the trivial bound $\|\nabla f(x_{k})\|^{2}\le (L\Delta_{\max})^{2}$
for the complement we obtain
\[
\min_{0\le i\le K}\|\nabla f(x_{i})\|
\;\le\;
\sqrt{\frac{2L\bigl(f(x_{0})-f^{*}\bigr)}{\eta_{1}\,|\mathcal{I}|}}.
\tag{24}
\]
Since $|\mathcal{I}|\ge K+1 - C$ with $C$ given by (20), for all sufficiently
large $K$ the term $C$ is negligible and we have
\[
\min_{0\le i\le K}\|\nabla f(x_{i})\|
\;\le\;
\frac{2L\bigl(f(x_{0})-f^{*}\bigr)}{\Delta_{\min}\sqrt{K+1}},
\]
where we used $\eta_{1}=0.1$ for simplicity of constants.
This is exactly the bound (8) asserted in Theorem \ref{thm:conv}.
∎

\item \textbf{Linear rate under strong convexity.}  
Assume additionally that $f$ is $\mu$‑strongly convex:
\[
f(y)\ge f(x)+\nabla f(x)^{\top}(y-x)+\frac{\mu}{2}\|y-x\|^{2}.
\tag{25}
\]
From (21) and the inequality $\|\nabla f(x_{k})\|\ge\mu\|x_{k}-x^{*}\|$ (a direct
consequence of (25)) we obtain
\[
f(x_{k+1})-f^{*}
\le\Bigl(1-\frac{\eta_{1}\mu}{L}\Bigr)\bigl(f(x_{k})-f^{*}\bigr).
\tag{26}
\]
Iterating (26) yields the linear convergence statement (9).
∎
\end{enumerate}

--------------------------------------------------------------------

\section*{Rate Derivation (With Constants)}
Collecting the constants appearing in the sublinear bound (8):
\[
C_{\text{sub}}=
\frac{2L\bigl(f(x_{0})-f^{*}\bigr)}{\Delta_{\min}} .
\]
Thus for any $K\ge1$,
\[
\boxed{\displaystyle
\min_{0\le k\le K}\|\nabla f(x_{k})\|
\le \frac{C_{\text{sub}}}{\sqrt{K+1}} } .
\]
If $f$ is $\mu$‑strongly convex, the linear factor in (9) is
\[
\rho:=1-\frac{\eta_{1}\mu}{L}\in(0,1),
\qquad
\boxed{\displaystyle
f(x_{K})-f^{*}\le\rho^{K}\bigl(f(x_{0})-f^{*}\bigr)} .
}
\]

--------------------------------------------------------------------

\section*{Special Cases}
\begin{itemize}
\item \textbf{Exact line search on quadratics.}  
When $f$ is a quadratic with Hessian $H\preceq L I$, the model (1) coincides
with $f$ itself; the trust‑region step becomes the exact Newton step
projected onto the ball.  In this case $\Delta_{k}$ can be kept constant,
and the algorithm converges in one iteration.

\item \textbf{Large radius regime.}  
If $\Delta_{k}\ge \|\nabla f(x_{k})\|/L$ for all $k$, the interior case (4) is always
selected and the method reduces to the classical gradient step with stepsize
$1/L$, reproducing the standard $O(1/k)$ convergence of gradient descent.

\item \textbf{Small radius regime.}  
When $\Delta_{k}$ is tiny, the algorithm behaves like a scaled subgradient
method with step length $\Delta_{k}$; the convergence rate remains
sublinear because the radius is never allowed to vanish (Assumption
\ref{ass:params}).
\end{itemize}

\bigskip
\noindent\textbf{End of document.}

\end{document}